\chapter{Lesson 11: The \ja{て}-Form and Its Uses}

\section{Why the \ja{て}-form matters}

The \ja{て}-form is a connecting form of the verb. By itself it does not show tense; it
links a verb to what follows. It is the building block for requests, permission,
prohibition, ongoing actions and joining sentences, and later for many more patterns.
It is the single most drilled form of first year.

\section{Forming the \ja{て}-form}

\subsection{\ja{る}-verbs and irregular verbs}
Easy: \ja{る}-verbs drop \ja{る} and add \ja{て}: \ja{たべる} $\to$ \ja{たべて},
\ja{みる} $\to$ \ja{みて}. Irregular: \ja{する} $\to$ \ja{して}, \ja{くる} $\to$ \ja{きて}.

\subsection{\ja{う}-verbs: depends on the last syllable}

\begin{center}
\begin{tikzpicture}[node distance=0.55cm and 0.6cm,
    start/.style={draw, rounded corners, fill=gray!15, font=\small, align=center, minimum height=0.8cm},
    q/.style={draw, diamond, aspect=2.2, fill=yellow!10, font=\small, align=center, inner sep=1pt},
    r/.style={draw, rounded corners, fill=green!10, font=\small, align=center, minimum width=3.1cm, minimum height=0.75cm}]
  \node[start] (s) at (0,0.0) {dictionary form};
  \node[q] (q0) at (0,-1.77) {\ja{する}/\ja{くる}?};
  \node[r, fill=red!10] (irr) at (5.2,-1.77) {\ja{して} / \ja{きて}};
  \node[q] (q1) at (0,-3.66) {\ja{る}-verb?};
  \node[r] (ru) at (5.2,-3.66) {drop \ja{る} + \ja{て}\\ \ja{たべる}$\to$\ja{たべて}};
  \node[q] (q2) at (0,-5.55) {\ja{いく}?};
  \node[r, fill=red!10] (iku) at (5.2,-5.55) {\ja{いって} (exception)};
  \node[q] (q3) at (0,-7.43) {ends in\\ \ja{う}・\ja{つ}・\ja{る}?};
  \node[r] (tte) at (5.2,-7.43) {\ja{って}\\ \ja{かう}$\to$\ja{かって}};
  \node[q] (q4) at (0,-9.32) {ends in\\ \ja{む}・\ja{ぶ}・\ja{ぬ}?};
  \node[r] (nde) at (5.2,-9.32) {\ja{んで}\\ \ja{のむ}$\to$\ja{のんで}};
  \node[q] (q5) at (0,-11.21) {ends in \ja{く}?};
  \node[r] (ite) at (5.2,-11.21) {\ja{いて}\\ \ja{かく}$\to$\ja{かいて}};
  \node[q] (q6) at (0,-13.1) {ends in \ja{ぐ}?};
  \node[r] (ide) at (5.2,-13.1) {\ja{いで}\\ \ja{およぐ}$\to$\ja{およいで}};
  \node[r] (shite) at (0,-14.87) {ends in \ja{す}: \ja{して}\\ \ja{はなす}$\to$\ja{はなして}};
  \draw[-Stealth] (s) -- (q0);
  \draw[-Stealth] (q0) -- node[above, font=\scriptsize] {yes} (irr);
  \draw[-Stealth] (q0) -- node[right, font=\scriptsize] {no} (q1);
  \draw[-Stealth] (q1) -- node[above, font=\scriptsize] {yes} (ru);
  \draw[-Stealth] (q1) -- node[right, font=\scriptsize] {no (\ja{う}-verb)} (q2);
  \draw[-Stealth] (q2) -- node[above, font=\scriptsize] {yes} (iku);
  \draw[-Stealth] (q2) -- node[right, font=\scriptsize] {no} (q3);
  \draw[-Stealth] (q3) -- node[above, font=\scriptsize] {yes} (tte);
  \draw[-Stealth] (q3) -- node[right, font=\scriptsize] {no} (q4);
  \draw[-Stealth] (q4) -- node[above, font=\scriptsize] {yes} (nde);
  \draw[-Stealth] (q4) -- node[right, font=\scriptsize] {no} (q5);
  \draw[-Stealth] (q5) -- node[above, font=\scriptsize] {yes} (ite);
  \draw[-Stealth] (q5) -- node[right, font=\scriptsize] {no} (q6);
  \draw[-Stealth] (q6) -- node[above, font=\scriptsize] {yes} (ide);
  \draw[-Stealth] (q6) -- node[right, font=\scriptsize] {no} (shite);
\end{tikzpicture}
\end{center}

\begin{center}
\begin{tabular}{llll}
\toprule
Ending & \ja{て}-form ending & Example & Meaning \\
\midrule
\ja{う}, \ja{つ}, \ja{る} & \ja{って} & \ja{あう}$\to$\ja{あって}, \ja{まつ}$\to$\ja{まって}, \ja{とる}$\to$\ja{とって} & meet, wait, take \\
\ja{む}, \ja{ぶ}, \ja{ぬ} & \ja{んで} & \ja{よむ}$\to$\ja{よんで}, \ja{あそぶ}$\to$\ja{あそんで}, \ja{しぬ}$\to$\ja{しんで} & read, play, die \\
\ja{く} & \ja{いて} & \ja{きく}$\to$\ja{きいて}, \ja{かく}$\to$\ja{かいて} & listen, write \\
\ja{ぐ} & \ja{いで} & \ja{いそぐ}$\to$\ja{いそいで} & hurry \\
\ja{す} & \ja{して} & \ja{はなす}$\to$\ja{はなして}, \ja{かす}$\to$\ja{かして} & speak, lend \\
exception & \ja{って} & \ja{いく}$\to$\ja{いって} & go \\
\bottomrule
\end{tabular}
\end{center}

\begin{aside}
Mnemonic (to the tune of ``Clementine''):
``\ja{う}, \ja{つ}, \ja{る} --- \ja{って}; \ja{む}, \ja{ぶ}, \ja{ぬ} --- \ja{んで};
\ja{く} --- \ja{いて}, \ja{ぐ} --- \ja{いで}; \ja{す} --- \ja{して}; \ja{いく} ---
\ja{いって}!''
\end{aside}

\begin{important}[Watch the exceptions]
\begin{itemize}
  \item \ja{いく} $\to$ \ja{いって} (not *\ja{いいて}).
  \item \ja{かえる}, \ja{はいる} are \ja{う}-verbs: \ja{かえって}, \ja{はいって}
        (not *\ja{かえて}).
  \item \ja{きる} ``wear'' (\ja{る}-verb) $\to$ \ja{きて}; \ja{きる} ``cut'' (\ja{う}-verb)
        $\to$ \ja{きって}.
\end{itemize}
\end{important}

\section{Use 1: requests --- \ja{〜てください}}
``Please do \ldots''
\begin{example}
\leavevmode
\begin{itemize}
  \item \ja{ここになまえをかいてください。} \emph{Koko ni namae o kaite kudasai.}
        Please write your name here.
  \item \ja{ちょっとまってください。} Please wait a moment.
  \item \ja{ペンをかしてください。} Please lend me a pen.
  \item \ja{ゆっくりはなしてください。} Please speak slowly.
\end{itemize}
\end{example}
\begin{remark}
\ja{〜てください} is polite but still a directive. To ask a favour of a teacher or
superior, \ja{〜てくださいませんか} (``would you please\ldots?'') is softer.
\end{remark}

\section{Use 2: permission --- \ja{〜てもいいです}}
``It's OK to \ldots / You may \ldots''. As a question, \ja{〜てもいいですか} asks
permission.
\begin{example}
\leavevmode
\begin{itemize}
  \item[Q:] \ja{しゃしんをとってもいいですか。} \emph{Shashin o totte mo ii desu ka.}
        May I take a picture?
  \item[A:] \ja{ええ、どうぞ。} Sure, go ahead.
  \item[A:] \ja{すみません、ちょっと\ldots} Sorry, (that's) a bit\ldots (= no)
  \item \ja{トイレにいってもいいですか。} May I go to the washroom?
  \item \ja{じしょをつかってもいいです。} You may use a dictionary.
\end{itemize}
\end{example}

\section{Use 3: prohibition --- \ja{〜てはいけません}}
``You must not \ldots'' Strong; used for rules, signs, and by people in authority.
\begin{example}
\leavevmode
\begin{itemize}
  \item \ja{ここでたばこをすってはいけません。} You must not smoke here.
  \item \ja{しけんのときにはなしてはいけません。} You must not talk during the exam.
  \item[Q:] \ja{ここでたべてもいいですか。} May I eat here?
  \item[A:] \ja{いいえ、たべてはいけません。} No, you may not.
\end{itemize}
\end{example}
In speech, \ja{〜ては} often contracts to \ja{〜ちゃ}: \ja{たべちゃいけない}.

\section{Use 4: ongoing action and state --- \ja{〜ている}}

\ja{〜ています} has two main meanings, depending on the verb:
\begin{center}
\begin{tikzpicture}[x=1cm,y=1cm]
  % action in progress
  \draw[-Stealth, thick] (0,0) -- (6,0) node[right, font=\small] {time};
  \fill[blue!20] (1.2,-0.25) rectangle (4.8,0.25);
  \draw[very thick, red!70!black] (3,-0.45) -- (3,0.45) node[above, font=\scriptsize] {now};
  \node[font=\small, anchor=east] at (-0.2,0) {\ja{よんでいます}};
  \node[font=\scriptsize] at (3,-0.75) {action in progress: ``is reading''};
  % resulting state
  \begin{scope}[yshift=-2.1cm]
  \draw[-Stealth, thick] (0,0) -- (6,0) node[right, font=\small] {time};
  \fill[red!70!black] (1.2,0) circle (0.12);
  \node[font=\scriptsize] at (1.2,0.45) {got married};
  \fill[green!20] (1.3,-0.25) rectangle (5.8,0.25);
  \draw[very thick, red!70!black] (3,-0.45) -- (3,0.45) node[above, font=\scriptsize] {now};
  \node[font=\small, anchor=east] at (-0.2,0) {\ja{けっこんしています}};
  \node[font=\scriptsize] at (3.5,-0.75) {resulting state: ``is married''};
  \end{scope}
\end{tikzpicture}
\end{center}

\begin{itemize}
  \item \textbf{Action in progress} (activity verbs): \ja{いまテレビをみています。}
        ``I'm watching TV now.'' \ja{あめがふっています。} ``It's raining.''
  \item \textbf{Resulting state} (change-of-state verbs such as \ja{けっこんする}
        marry, \ja{しる} know, \ja{すむ} live, \ja{ふとる} gain weight, \ja{もつ} hold):
        \ja{あねはけっこんしています。} ``My older sister is married.''
        \ja{トロントにすんでいます。} ``(I) live in Toronto.''
  \item \textbf{Habitual} actions: \ja{まいにちジムにいっています。}
        ``I go to the gym every day (these days).''
\end{itemize}

\begin{important}[\ja{しっています} vs.\ \ja{しりません}]
``I know'' is \ja{しっています}, but ``I don't know'' is \ja{しりません}, not
*\ja{しっていません}. Also note \ja{いっています} can mean ``has gone (and is there
now)'', not ``is going''.
\end{important}

\section{Use 5: joining actions --- \ja{〜て、〜}}
The \ja{て}-form links verbs in sequence (``and then''). Tense comes from the final verb.
\begin{example}
\ja{あさおきて、シャワーをあびて、だいがくにいきました。}
\emph{Asa okite, shawā o abite, daigaku ni ikimashita.}
I got up in the morning, showered, and went to the university.\\
\ja{としょかんにいって、べんきょうします。} I'll go to the library and study.
\end{example}
\ja{い}-adjectives and \ja{な}-adjectives/nouns have \ja{て}-forms too:
\ja{やすくて、おいしい} ``cheap and tasty''; \ja{しずかで、きれい} ``quiet and
pretty''.

\section{Kanji}
\begin{center}
\begin{tabular}{lllll}
\toprule
Kanji & Meaning & Readings & Example & \\
\midrule
\ja{行} & go & \ja{コウ}; \ja{い(く)} & \ja{行って} & itte \\
\ja{来} & come & \ja{ライ}; \ja{く(る)}, \ja{き(ます)} & \ja{来週} & raishū, next week \\
\ja{見} & see & \ja{ケン}; \ja{み(る)} & \ja{見てください} & mite kudasai \\
\ja{食} & eat & \ja{ショク}; \ja{た(べる)} & \ja{食べもの} & tabemono, food \\
\ja{飲} & drink & \ja{イン}; \ja{の(む)} & \ja{飲みもの} & nomimono, drink \\
\ja{聞} & hear, ask & \ja{ブン}; \ja{き(く)} & \ja{新聞} & shinbun, newspaper \\
\ja{読} & read & \ja{ドク}; \ja{よ(む)} & \ja{読んで} & yonde \\
\ja{書} & write & \ja{ショ}; \ja{か(く)} & \ja{書いて} & kaite \\
\bottomrule
\end{tabular}
\end{center}

\begin{center}
\begin{tikzpicture}[k/.style={draw, minimum size=1.2cm, font=\huge}]
  \node[k] (a) at (0,0) {\ja{門}};
  \node[font=\Large] at (1.1,0) {+};
  \node[k] (b) at (2.2,0) {\ja{耳}};
  \node[font=\Large] at (3.3,0) {=};
  \node[k, fill=blue!8] (c) at (4.4,0) {\ja{聞}};
  \node[font=\small] at (0,-0.95) {gate};
  \node[font=\small] at (2.2,-0.95) {ear};
  \node[font=\small] at (4.4,-0.95) {listen};
\end{tikzpicture}
\end{center}
An ear at the gate: listening. Many verb kanji have a meaning component plus a sound
component; \ja{飲} and \ja{食} share the ``food'' radical \ja{飠}.

\section*{Lesson 11 review}
\begin{itemize}
  \item \ja{る}-verbs: \ja{〜て}. \ja{う}-verbs: \ja{うつる}$\to$\ja{って},
        \ja{むぶぬ}$\to$\ja{んで}, \ja{く}$\to$\ja{いて}, \ja{ぐ}$\to$\ja{いで},
        \ja{す}$\to$\ja{して}; \ja{いく}$\to$\ja{いって}. \ja{して}, \ja{きて}.
  \item \ja{〜てください} request; \ja{〜てもいいですか} permission;
        \ja{〜てはいけません} prohibition.
  \item \ja{〜ています}: in progress, resulting state, or habit.
  \item \ja{て} joins actions in sequence.
  \item Kanji: \ja{行来見食飲聞読書}.
\end{itemize}
